\documentclass[12pt,a4paper]{article}
\usepackage{physics-textbook}

\begin{document}
\BookParagraph{0}{Основы решения физических задач}
Физическую задачу начинают с описания ситуации. Нужно установить, какие тела
участвуют в движении, какие величины известны и что требуется определить.
После этого выбирают физический закон и записывают его для рассматриваемого
случая.

\BookSubsection{0.1}{От ситуации к уравнению}
При равноускоренном прямолинейном движении изменение скорости
пропорционально времени. Если начальная скорость равна нулю, то
\[
v=at,
\]
где $v$ — скорость (м/с), $a$ — ускорение (м/с$^2$), $t$ — время (с).

\begin{figure}[ht]
\centering
\begin{tikzpicture}[book diagram]
  \draw (0,0)--(6,0);
  \draw[book body] (1,0) rectangle (1.8,.5);
  \draw[book force] (1.8,.25)--(3,.25) node[above] {$\vec F$};
  \draw[book displacement] (1,-.5)--(4,-.5) node[midway,below] {$\vec s$};
\end{tikzpicture}
\caption{Сила и перемещение тела}
\end{figure}

Сначала получают выражение для искомой величины, затем подставляют числовые
значения с согласованными единицами измерения.
\end{document}
